\documentclass[10pt,a4paper]{amsart}
\usepackage[margin=19mm]{geometry}
\usepackage{amsmath,amssymb,mathtools}
\usepackage[scr=rsfs,cal=euler]{mathalfa}
\usepackage{xfrac}
\usepackage[unicode,hidelinks,bookmarks=false]{hyperref}
\newcommand{\ie}{i.e.\ }
\newcommand{\cf}{cf.\ }
\newcommand{\resp}{respectively\ }
\newcommand{\Subsection}{\subsection}
\providecommand{\nobreakdash}{\nobreak\hbox{-}}
\setlength{\emergencystretch}{2em}
\multlinegap0pt
\usepackage{amssymb}
\usepackage{enumerate}
\usepackage{xcolor}
\usepackage{bm}
\usepackage[shortlabels]{enumitem}
\usepackage{dsfont}
\newtheorem{lemma}{Lemma}
\newtheorem{theorem}{Theorem}
\title{Embedding martingale diffusions as binary posteriors in sequential inference}
\author{Steven Campbell$^a$ and Karl Kristian Engelund$^b$}
\date{Source: arXiv:2607.20373v1 (CC BY 4.0)}
\begin{document}
\maketitle

\noindent\emph{Source and licence.} Embedding martingale diffusions as binary posteriors in sequential inference. Version arXiv:2607.20373v1 (CC BY 4.0), distributed by arXiv under the Creative Commons Attribution 4.0 International licence, http://creativecommons.org/licenses/by/4.0/, as recorded in the arXiv licence metadata for this exact version and shown on the abstract page https://arxiv.org/abs/2607.20373v1. Full licence notice: \texttt{licenses/arxiv-2607.20373-CC-BY-4.0.txt}.\par\smallskip

\noindent\textit{Editorial scope.} Counted unit: the article's theorem thm:posterior-embedding (source0.tex lines 743--755). The article's complete original proof of it is delivered with it (source0.tex lines 757--802, one \texttt{proof} environment of 46 lines). No mathematics is written, repaired or re-derived: the statement and the proof are the article's own lines, in the article's own notation, and no retained line is substituted or re-flowed.  Retained uncounted as the source-local context the proof needs: lines 168--171; lines 226--226; lines 235--240; lines 243--246; lines 359--365; lines 540--542; lines 624--638; lines 726--730.  Nothing was omitted from any retained span, so the package carries no omission record.  No retained line contains a citation, so the wrapper carries no bibliography and no bibliography entry.  No retained line contains a figure environment, an image-inclusion command or drawing code, so no image is delivered and the package declares no asset dependency.  Editorial formatting only, in addition: an AMS \texttt{amsart} preamble that restates the article's own declarations and macros used by the retained lines, namely the environments \texttt{lemma}, \texttt{theorem} on separate counters, together with the standard packages those lines need.  The retained environments print in this document as Lemma 1, Theorem 1 in order of appearance, whereas the article prints the same items with the numbering of its own full document.  The complete native source of the article as distributed in the arXiv e-print for this exact version is bundled unchanged as \texttt{sources/205807/source0.tex}.
\begin{equation}\label{eq:intro-target}
\Pi_t=\pi+\int_0^{t\wedge\tau}\sigma(\Pi_s)dB_s,
  \qquad t\geq0, \qquad \pi\in(0,1).
\end{equation}

\subsection{Binary sequential inference}\label{sec:bin-seq-inf}

\begin{equation}\label{eq:signal-process}
  X_t=\begin{cases}
      x_0 + \int_0^{t}\mu_\theta(X_s)ds+W_{t}, & t<\zeta\\
      X_\zeta, & t\geq \zeta
  \end{cases}, \qquad t\geq0,
\end{equation}

\begin{lemma}\label{lem:signal-well-posed}
Equation
\eqref{eq:signal-process} admits a pathwise unique strong solution.
\end{lemma}

\begin{equation}\label{eq:filtering-equation}
  P_t
  =
  p+\int_0^t
    P_s(1-P_s)\delta(X_s)\,d\overline W_s,
  \qquad t<\zeta.
\end{equation}

\begin{equation}\label{eq:F-def}
  F(p)=\int_{1/2}^p \frac{1}{\sigma(u)}du+C, \qquad p\in(0,1),
\end{equation}

\begin{equation}\label{eq:embedding-densities}
  Z_t^1
  :=
  \frac{d\mathbb Q^1}{d\mathbb P}
  \bigg|_{\mathcal F_t}
  =
  \frac{\Pi_t}{\pi},
  \qquad
  Z_t^0
  :=
  \frac{d\mathbb Q^0}{d\mathbb P}
  \bigg|_{\mathcal F_t}
  =
  \frac{1-\Pi_t}{1-\pi}.
\end{equation}

\begin{equation}\label{eq:embedding-drifts}
  \nu_1(y)=-\frac12\sigma'(G(y))+\frac{\sigma(G(y))}{G(y)},
  \qquad
  \nu_0(y)=-\frac12\sigma'(G(y))-\frac{\sigma(G(y))}{1-G(y)}.
\end{equation}

\begin{theorem}[Posterior embedding]
\label{thm:posterior-embedding}
Let $\Pi$ be the autonomous win-martingale in
\eqref{eq:intro-target}, and let $F$, $G$, $I$, $\nu_0$, and $\nu_1$
be defined by \eqref{eq:F-def} and \eqref{eq:embedding-drifts}.
In the binary sequential experiment of
Section \ref{sec:bin-seq-inf}, take $p=\pi$, $J=I$, $x_0=F(\pi)$, and $\mu_i\equiv\nu_i$, $i\in\{0,1\}$.
Then, $P=G(X)$ up to indistinguishability,
$\zeta=\inf\{t\geq0:P_t\in\{0,1\}\}$,
\begin{equation}\label{eq:thm-P-embedding}
  P_t=p+\int_0^{t\wedge \zeta}\sigma(P_s)\,d\overline W_s, \quad t\geq0, \quad \text{and} \quad \operatorname{Law}_{\widetilde{\mathbb P}}(P)=\operatorname{Law}_{\mathbb P}(\Pi).
\end{equation}
\end{theorem}

\begin{proof}
By Lemma \ref{lem:signal-well-posed}, the prescribed experiment is
well defined. Recall that $\widetilde{\mathbb Q}^i:=\widetilde{\mathbb P}(\,\cdot\mid\theta=i)$ for $i\in\{0,1\}$. Under $\widetilde{\mathbb Q}^i$, the signal $X$ solves the absorbed equation $dX_t=\nu_i(X_t)\,dt+dW_t$ for $t<\zeta$.
At the same time, the preceding Girsanov calculation shows that $Y=F(\Pi)$ satisfies
the same equation under $\mathbb Q^i$. Uniqueness in law therefore gives $\operatorname{Law}_{\widetilde{\mathbb Q}^i}(X)=\operatorname{Law}_{\mathbb Q^i}(Y)$, $i\in\{0,1\}$. Since $p=\pi$ and $\widetilde{\mathbb P}
  =
  (1-\pi)\widetilde{\mathbb Q}^0
  +
  \pi\widetilde{\mathbb Q}^1$, $\mathbb P
  =
  (1-\pi)\mathbb Q^0+\pi\mathbb Q^1$,
it follows that $\operatorname{Law}_{\widetilde{\mathbb P}}(X)=\operatorname{Law}_{\mathbb P}(Y)$.

Fix $t\geq0$ and let $\Phi$ be any bounded Borel functional of the signal path up to time $t$. Since $\Pi=G(Y)$, we have by the density
identities \eqref{eq:embedding-densities}
\[
  \widetilde{\mathbb E}\bigl[\theta\Phi(X)\bigr]
  =
  \pi\,
  \mathbb E_{\widetilde{\mathbb Q}^1}\bigl[\Phi(X)\bigr] =
  \pi\,
  \mathbb E_{\mathbb Q^1}\bigl[\Phi(Y)\bigr] =
  \mathbb E\bigl[\Pi_t\Phi(Y)\bigr] =
  \widetilde{\mathbb E}
  \bigl[G(X_t)\Phi(X)\bigr].
\]
It follows that $P_t= \widetilde{\mathbb E}[\theta\mid\mathcal F_t^X]=G(X_t)$ $\widetilde{\mathbb P}$--a.s.
Equality at rational times, together with the c\`adl\`ag property of
$P$ and continuity of $G(X)$, gives indistinguishability. Since $G(I)=(0,1)$ and $G$ maps the two endpoints of $\overline I$
to $0$ and $1$, respectively, $\zeta
  =
  \inf\{t\geq0:P_t\in\{0,1\}\}$,
and $P=G(X)$ is absorbed after $\zeta$. Finally, the definitions of $\nu_0$ and $\nu_1$ give $\delta(x):=
  \nu_1(x)-\nu_0(x) 
  %
  %
  %
  =
  \sigma(G(x))/(G(x)(1-G(x)))$.
Using $P_t=G(X_t)$ in the filtering equation
\eqref{eq:filtering-equation}, we obtain $dP_t
  =
  P_t(1-P_t)\delta(X_t)\,d\overline W_t=
  \sigma(P_t)\,d\overline W_t$, $t<\zeta$,
and \eqref{eq:thm-P-embedding} follows by absorption. The equality of laws is immediate.
\end{proof}

\end{document}
